% Abstract: tight, claim-light. Numbers go in only after the v1
% experiment runs. Until then, keep methodological framing.
Hyperparameter tuning of gradient-boosted decision trees is inherently
multi-objective: predictive quality must be balanced against
computational cost and, in practice, against robustness across data
splits. We present a design-aware framework that integrates a Design of
Experiments (DoE) phase, design-aware response surface modeling
(process quadratic on factorial and central composite designs;
Scheff\'{e} canonical mixture polynomials on simplex weight designs), a
flexible factor-analysis layer (PCA with Varimax rotation, configurable
between automatic, fixed, manual-construct, and disabled modes), and a
true \emph{$N$-objective} Normal Boundary Intersection (NBI) routine
implementing the Das--Dennis sub-problem with explicit anchors, payoff
matrix, quasi-normal direction, and equality constraints rather than a
weighted-sum scalarization. A conditional Mixture-Based Performance
Assessment (MBPA) post-optimization stage refines weights when the
primary frontier is flat or weakly informative. We instantiate the
framework on the
$12{\times}3{\times}10$ benchmark introduced in
Section~\ref{sec:experimental-design} (twelve public tabular datasets,
three GBDT families, ten replicas under matched evaluation budgets) and
position it as a methodological refinement and generalization of the
dissertation pipeline of \citep{ribeiro2026dissertation}, anchored in
the multivariate NBI/MBPA formulation of \citep{pereira2025eaai}.
\todoblock{Insert quantitative claims (median CPU-hour reduction,
predictive parity bands, MBPA trigger rates) once the v1 experiment is
aggregated; write each claim verbatim from the result CSVs to avoid
drift.}
